\section{Direcciones futuras}

Shunyu Yao propuso cinco direcciones de investigación clave:

\subsection{Entrenamiento (Training)}

\begin{importantbox}{Coevolución de modelo y Agent}
Los datos de entrenamiento de los LLM actuales provienen principalmente de texto de internet y carecen de datos de trayectorias de Agent (proceso de razonamiento, secuencia de acciones, retroalimentación del entorno). Solución: usar un prompted agent para generar grandes cantidades de datos de trayectorias, y luego usar esos datos para hacer fine-tuning al modelo, formando un ciclo positivo. Es análogo a la coevolución entre las GPU y el aprendizaje profundo.
\end{importantbox}

\subsection{Interfaz (Interface)}

\begin{itemize}
    \item La interfaz humano-máquina (HCI) se ha estudiado durante décadas, pero la interfaz óptima para un Agent es distinta de la que conviene a los humanos
    \item Por ejemplo: los humanos buscan archivos con \texttt{ls} + \texttt{cd} (un resultado a la vez), pero a un Agent le conviene más obtener varios resultados de una vez
    \item No se trata de optimizar al Agent, sino de optimizar el \textbf{entorno}: el mismo modelo se desempeña mejor con una mejor interfaz
\end{itemize}

\subsection{Robustez (Robustness)}

\begin{warningbox}{Pass@K vs. robustez}
El ámbito académico se enfoca en pass@K (tener éxito al menos una vez en K intentos), pero las aplicaciones reales necesitan \textbf{robustez} (tener éxito en cada uno de los K intentos). En todos los modelos actuales, la robustez disminuye de manera monótona a medida que aumenta el número de muestreos. Un Agent de servicio al cliente puede perder un cliente con un solo fallo: no basta con aspirar a «poder lograrlo», también hay que «poder lograrlo siempre».
\end{warningbox}

\subsection{Colaboración humano-máquina (Human-in-the-loop)}

En escenarios reales, el Agent necesita interactuar con personas: el cliente no da toda la información de una sola vez, y el Agent necesita varias rondas de conversación para obtener la información que requiere.

\subsection{Bancos de pruebas (Benchmarks)}

\begin{itemize}
    \item Se necesitan bancos de pruebas más realistas y de mayor escala (como Tau-Bench: escenarios de servicio al cliente)
    \item Se necesita evaluar la robustez y no solo el límite superior de la capacidad
    \item Se necesita incluir elementos de interacción humano-máquina
\end{itemize}

\subsection{Resumen de la sección}

Las cinco direcciones clave del campo de los Agent de LLM (entrenamiento, interfaz, robustez, colaboración humano-máquina y evaluación) tienen mucha fruta fácil de alcanzar y no requieren cantidades enormes de recursos de GPU, por lo que son adecuadas para la investigación académica.

